\section{Paired-intervention identifiability calculations}
This section is a theoretical and synthetic extension motivating the next Paired-Acquisition Neural Factorization experiments. It is separate from the detector-switch evidence above and should not be interpreted as a clinical validation result.

\subsection{Why pairing changes the information structure}
Assume a pathology representation is generated by biological and acquisition factors:
\begin{equation}
  x=Bz_b+Az_a+\varepsilon.
\end{equation}
For a paired acquisition intervention preserving biology,
\begin{align}
  x_i&=Bz_{b,i}+Az_{a,i}+\varepsilon_i,\\
  x_i'&=Bz_{b,i}+Az'_{a,i}+\varepsilon_i'.
\end{align}
Subtracting gives
\begin{equation}
  \Delta x_i=x_i'-x_i=A(z'_{a,i}-z_{a,i})+(\varepsilon_i'-\varepsilon_i),
\end{equation}
so the shared biological term cancels. If acquisition latents have covariance $\Sigma_a$ and the noise is isotropic,
\begin{equation}
  \operatorname{Cov}(\Delta x)=2A\Sigma_aA^\top+2\sigma^2I.
\end{equation}
Thus the leading eigenspace of paired differences estimates the acquisition subspace. A first-order biological projection is
\begin{equation}
  P_{\perp A}=I-AA^\top,\qquad x_{\mathrm{bio}}=P_{\perp A}x.
\end{equation}

\subsection{Geometric overlap penalty}
Let $B$ and $A$ be orthonormal bases with canonical overlap $\rho$. Removing the acquisition subspace retains approximately
\begin{equation}
  E_{\mathrm{bio,retained}}\approx1-\rho^2.
\end{equation}
Hence $\rho=0.90$ leaves only $0.19$ of the biological energy outside the acquisition subspace, while $\rho=0.95$ leaves $0.0975$. Any component in $\operatorname{span}(B)\cap\operatorname{span}(A)$ is not uniquely assignable without additional assumptions.

\begin{table}[t]
\centering
\caption{Geometric difficulty proxy.}
\small
\begin{tabular}{ccc}
\toprule
$\rho$ & $1-\rho^2$ & $(1-\rho^2)^{-2}$\\
\midrule
0.50 & 0.7500 & 1.78\\
0.75 & 0.4375 & 5.22\\
0.90 & 0.1900 & 27.70\\
0.95 & 0.0975 & 105.19\\
\bottomrule
\end{tabular}
\end{table}

\subsection{Finite-pair sample complexity}
Let $m=pn$ be the number of paired examples, $D$ the observed dimension, and $\gamma$ the eigengap separating acquisition directions from noise. A Davis--Kahan-style perturbation argument suggests
\begin{equation}
  \sin\Theta(\widehat A,A)
  \lesssim
  \frac{C}{\gamma}
  \sqrt{\frac{D+\log(1/\delta)}{pn}}.
\end{equation}
Requiring this estimation error to be smaller than the surviving biological margin gives the provisional scaling prediction
\begin{equation}
  p_{\min}
  \gtrsim
  \frac{C^2(D+\log(1/\delta))}
  {n\gamma^2(1-\rho^2)^2}.
\end{equation}
This predicts diminishing returns in paired coverage, error decay proportional to $(pn)^{-1/2}$, and a sharp rise in required pairing as $\rho\rightarrow1$. The exact exponent and overlap dependence remain hypotheses rather than established theorems for the learned nonlinear models.

\subsection{Synthetic two-resource recovery boundary}
A follow-up synthetic phase-map study treated unpaired observations and paired interventions as distinct resources. It varied observational sample size $N_u\in\{500,1000,3000,10000\}$, high biological--acquisition overlap $\rho\in\{0.90,0.95\}$, linear versus nonlinear mixing, a grid of paired-anchor counts up to 300, five seeds, and pairing-aware methods. The endpoint $N_p^{\ast}$ was the smallest paired-anchor count at which the mean universal biological score crossed a predefined threshold and remained above it for every larger tested pair count.

At the primary sustained-recovery threshold $\tau=0.50$, 43 of 47 available conditions had observed boundaries and four were right-censored above the tested pair-count range. A model using $\log N_u$, overlap, and nonlinear mixing achieved the best leave-condition-out prediction and slightly outperformed a larger model with method-specific offsets.

\begin{table}[t]
\centering
\caption{Recovery-boundary model comparison at $\tau=0.50$. Lower leave-condition-out log RMSE is better.}
\small
\begin{tabular}{lc}
\toprule
Model & Log RMSE\\
\midrule
Power + overlap + nonlinearity & 0.455\\
Full model with method offsets & 0.459\\
Power only & 0.599\\
Constant boundary & 0.881\\
\bottomrule
\end{tabular}
\end{table}

The uncensored empirical fit was
\begin{align}
  \log N_p^{\ast}
  &=-1.568-0.659\log N_u \\
  &\quad+10.8792\rho+0.5345V,
\end{align}
where $V=1$ denotes nonlinear mixing and $V=0$ denotes linear mixing. Equivalently,
\begin{align}
  N_p^{\ast}
  &\approx0.208\,N_u^{-0.659} \\
  &\quad\times\exp(10.8792\rho)\exp(0.5345V).
\end{align}

An interval/right-censored likelihood fit used all 47 conditions and treated the four missing primary boundaries as exceeding the maximum tested pair count. It yielded
\begin{align}
  \log N_p^{\ast}
  &=-5.077-0.8086\log N_u \\
  &\quad+15.8347\rho+0.5277V,
\end{align}
and therefore
\begin{align}
  N_p^{\ast}
  &\approx0.00624\,N_u^{-0.8086} \\
  &\quad\times\exp(15.8347\rho)\exp(0.5277V).
\end{align}
The nonlinear multiplier is $\exp(0.5277)\approx1.695$. Within this tested synthetic regime, more observational data was associated with fewer required paired anchors, whereas overlap and nonlinear mixing increased the paired-anchor requirement.

The qualitative sign pattern was stable in ordinary least-squares sensitivity analyses at thresholds $0.45$, $0.50$, $0.55$, and $0.60$: the coefficient of $\log N_u$ remained negative, while the coefficients of overlap and nonlinear mixing remained positive. The number of observed boundaries decreased from 47 at $\tau=0.45$ to 27 at $\tau=0.60$, so stricter-threshold estimates are increasingly influenced by right censoring. Censored fits converged at $\tau=0.50$ and $\tau=0.55$; non-converged fits are not used as headline estimates.

These results support a threshold-dependent empirical family
\begin{equation}
  N_p^{\ast}(\tau)
  =C_{\tau}N_u^{-\alpha_{\tau}}
  \exp(\gamma_{\tau}\rho)
  \exp(\eta_{\tau}V),
\end{equation}
not a universal coefficient set. The raw-overlap term is a local approximation over the tested high-overlap range and does not by itself confirm the theoretical $(1-\rho^2)^{-2}$ dependence.
